\documentclass[11pt]{amsart}
\usepackage[T1]{fontenc}
\usepackage{lmodern,amsmath,amssymb,mathtools,booktabs,array,microtype}
\usepackage[margin=1in]{geometry}
\usepackage[hidelinks]{hyperref}
\usepackage{enumitem}
\usepackage[section]{placeins}
\setlist[enumerate]{leftmargin=*,itemsep=2pt,topsep=4pt}
\newtheorem{theorem}{Theorem}[section]
\newtheorem{proposition}[theorem]{Proposition}
\newtheorem{lemma}[theorem]{Lemma}
\newtheorem{corollary}[theorem]{Corollary}
\theoremstyle{definition}
\newtheorem{remark}[theorem]{Remark}
\newtheorem{example}[theorem]{Example}
\DeclareMathOperator{\FPdim}{FPdim}
\DeclareMathOperator{\rank}{rank}
\DeclareMathOperator{\Rep}{Rep}
\DeclareMathOperator{\Tr}{Tr}
\DeclareMathOperator{\Irr}{Irr}
\DeclareMathOperator{\Ad}{Ad}
\newcommand{\C}{\mathbb C}
\newcommand{\Q}{\mathbb Q}
\newcommand{\Z}{\mathbb Z}
\newcommand{\one}{\mathbf1}
\newcommand{\Fib}{\mathrm{Fib}}
\newcommand{\Near}{\mathrm{NG}}
\numberwithin{equation}{section}
\allowdisplaybreaks
\hypersetup{pdftitle={Classifying Grothendieck rings of pivotal fusion categories of rank four},pdfauthor={Jingcheng Dong, Sebastien Palcoux, and Arnaud Plessis}}

\title[Pivotal Grothendieck rings of rank four]{Classifying Grothendieck rings of pivotal fusion categories of rank four}
\author{Jingcheng Dong}
\address{College of Mathematics and Statistics, and Center for Applied Mathematics of Jiangsu Province, Nanjing University of Information Science and Technology, Nanjing 210044, China}
\email{jcdong@nuist.edu.cn}
\author{S\'ebastien Palcoux}
\address{Beijing Institute of Mathematical Sciences and Applications, Huairou District, Beijing 101408, China}
\email{sebastienpalcoux@gmail.com}
\author{Arnaud Plessis}
\address{Beijing Institute of Mathematical Sciences and Applications, Huairou District, Beijing 101408, China}
\email{plessis@bimsa.cn}
\subjclass[2020]{Primary 18M20; Secondary 16Z05, 18M15}
\keywords{fusion category, Grothendieck ring, Frobenius--Schur indicator, central induction}
\date{}
\begin{document}
\begin{abstract}
We classify the Grothendieck rings of pivotal fusion categories of rank four over the complex field. There are exactly fifteen, each admitting a unitary categorification. 
This classification is obtained in two stages. 
The first consists in proving that the global Frobenius--Perron dimension of the Grothendieck ring of any pivotal fusion category of rank four is less than $3600$. 
The second uses known and new exclusion criteria to show that, among all fusion rings of rank four with global Frobenius--Perron dimension less than $3600$, only fifteen admit a pivotal categorification.
\end{abstract}
\maketitle
\setcounter{tocdepth}{2}
\tableofcontents

\section{Introduction}
Fusion categories over the complex field connect representation theory
and subfactor theory with low-dimensional topology and mathematical
physics. Spherical fusion categories furnish state-sum invariants of
three-manifolds \cite{BarrettWestbury} and unitary fusion categories
provide the algebraic input for string-net models of topological phases
\cite{LevinWen}. Their classification has a strong arithmetic aspect:
dimensions and formal codegrees are algebraic integers constrained by
Galois theory and divisibility \cite{ENO,OstrikCodegrees}.
We refer to \cite{EGNO} for the foundations and standard conventions of
fusion category theory.

The \emph{rank} of a fusion category counts the isomorphism classes of
its simple objects; its \emph{multiplicity} is the largest coefficient
in the decomposition of a tensor product of two simple objects.
Fixing the rank alone does not bound the multiplicity of an abstract
fusion ring: the rank-two rings with $x^2=1+ax$ exist for every
$a\in\Z_{\geq0}$. A finite classification of Grothendieck rings of fixed rank therefore requires
constraints coming from categorifiability.
In rank four, we obtain, thanks to the interaction between arithmetic constraints, the characters of Grothendieck rings, and the Frobenius-Schur indicators, that the global Frobenius--Perron dimension of the Grothendieck ring of any pivotal fusion category of rank four is less than $3600$.  
It leads to the following classification.

\begin{theorem}\label{thm:main}
Up to fusion ring isomorphism, there are fifteen fusion rings of rank four that admit a pivotal categorification; they are listed in Table~\ref{tab:models}. Each of them admits a unitary categorification.
\end{theorem}

The theorem classifies fusion rings, with all categorifications of the
same fusion ring counted once.
It also establishes that the existence of a
pivotal categorification is equivalent, in rank four, to the existence of
a unitary categorification. 

\begin{table}[ht]
\centering\small
\begin{tabular}{@{}clcrl@{}}
\toprule
$\mu(R)$ & A unitary realization & $\FPdim(\mathcal C)$ & Approx. & Reference \\
\midrule
$1$ & $\mathrm{Vec}_{C_2\times C_2}$ & $4$ & $4.000$ & \cite[\S3.4]{LPR} \\
$1$ & $\mathrm{Vec}_{C_4}$ & $4$ & $4.000$ & \cite[\S3.4]{LPR} \\
$1$ & Tambara--Yamagami on $C_3$ & $6$ & $6.000$ & \cite{TambaraYamagami} \\
$1$ & $\Fib\boxtimes\mathrm{Vec}_{C_2}$ & $5+\sqrt{5}$ & $7.236$ & \cite[\S3.4]{LPR} \\
$1$ & $\Rep(D_5)$ & $10$ & $10.000$ & \cite[\S3.4]{LPR} \\
$1$ & $\Fib\boxtimes\Fib$ & $(15+5\sqrt{5})/2$ & $13.090$ & \cite[\S3.4]{LPR} \\
$1$ & $\mathrm{PSU}(2)_6$ & $8+4\sqrt{2}$ & $13.657$ & \cite[Theorem 3.1]{EIP} \\
$1$ & Principal even part of $S'$ & $8+4\sqrt{2}$ & $13.657$ & \cite[Theorem 1]{LMP} \\
$1$ & $\mathrm{PSU}(2)_7$ & $\delta_7$ & $19.234$ & \cite[\S3.4]{LPR} \\
$2$ & $\Rep(A_4)$ & $12$ & $12.000$ & \cite[Theorem 1.1]{Larson} \\
$2$ & $\mathcal H_1$ (Haagerup) & $(39+9\sqrt{13})/2$ & $35.725$ & \cite[Table 2]{GrossmanSnyder} \\
$2$ & Principal even part of $2D2$ & $20+8\sqrt{5}$ & $37.889$ & \cite[Theorem 3.1]{EIP} \\
$3$ & Near-group type $(C_3,3)$ & $(21+3\sqrt{21})/2$ & $17.374$ & \cite[Example 9.4]{IzumiNear} \\
$3$ & $\Ad(E_8)$ & $\delta_E$ & $35.238$ & \cite{IzumiE8,EvansJones} \\
$6$ & Near-group type $(C_3,6)$ & $24+12\sqrt{3}$ & $44.785$ & \cite[Theorem 10.19]{IzumiNear} \\
\bottomrule
\end{tabular}
\caption{One unitary realization of each surviving fusion ring, with its
multiplicity and global Frobenius--Perron dimension
(rounded to three decimals in the fourth column). Rows are ordered first
by multiplicity and then by increasing global Frobenius--Perron dimension.
The numbers $\delta_7$ and $\delta_E$ are the largest real roots of
$x^3-27x^2+162x-243$ and
$x^4-45x^3+375x^2-1125x+1125$, respectively. 
The table also records references for the construction of each of these fifteen fusion rings.}
\label{tab:models}
\end{table}

\begin{remark}\label{rem:sharp}
For pivotal fusion categories of rank four, Table~\ref{tab:models} gives the sharp bounds $\mu(K_0(\mathcal C))\leq6$ and $\FPdim(\mathcal C)\leq24+12\sqrt3$, both attained by the near-group fusion ring of type $(C_3,6)$.
\end{remark}

From now on, all fusion categories are over $\C$. 
Ostrik classified fusion categories of rank two \cite{OstrikRank2}
and pivotal fusion categories of rank three \cite{OstrikRank3}.
In rank four, the modular and premodular cases were studied in
\cite{RSW,HongRowell,Bruillard4}. Without braiding, Larson
\cite{Larson} reduced the non-self-dual pseudo-unitary case to seven
fusion rings. Dong--Zhang--Dai \cite{DongZhangDai} treated the
nontrivially graded self-dual spherical case, and
Edie-Michell--Izumi--Penneys \cite{EIP} classified
$\Z/2\Z$-quadratic unitary fusion categories. In particular, the ring
$S(2)$, already present as $K_2(2)$ in Larson's list
\cite[Theorem~1.1(6)]{Larson}, was excluded from unitary categorification. 
Example \ref{ex:S(2)} explains how a short calculation allows to rule out $S(2)$ from a pivotal categorification. 
Further classification results appear in \cite{DongChenWang,DongSelfDual4}. 

We now conclude this introduction by outlining the organization of the paper. 
Section \ref{sec:spherical} sets out the basic terminology for fusion rings and fusion categories, then recalls several well-known results on categorifiable fusion rings. 
Section \ref{sec:central} establishes new inequalities, based on the Frobenius--Schur indicators, that must be satisfied by any pseudo-unitary fusion category of a given rank. 
Section \ref{sec:uniform} contains the crucial statement: the global Frobenius--Perron dimension of any pivotal fusion category of rank four is less than $3600$. 
The non-self-dual case follows from Larson \cite{Larson}, the self-dual pivotal, but not pseudo-unitary, case is handled by an arithmetic argument, and the remaining case involves inequalities coming from Section \ref{sec:central}.
The final Section \ref{sec:established} provides Python commands to find all fusion rings of rank four whose global Frobenius--Perron dimension is less than $3600$; there are $8977$. 
However, with the help of other commands, we will see that $8962$ of them violate at least one of the properties of pivotal fusion categories stated in Sections \ref{sec:spherical} and \ref{sec:central}. Theorem \ref{thm:main} then follows. 

\section{Preliminaries on fusion rings and fusion categories}\label{sec:spherical}

\subsection{Fusion rings and formal codegrees}
A \textit{fusion ring} $R$ is a free $\Z$-module of finite rank with basis $\{b_0=1, b_1,\dots, b_{r-1}\}$ and multiplication \[
 b_ib_j=\sum_kN_{ij}^kb_k,\] where $N_{ij}^k\geq 0$ are integers, called the \textit{fusion coefficients}, satisfying:  
\begin{itemize}
\item (Associativity) $\forall i,j,k$, one has $b_i(b_jb_k)=(b_ib_j)b_k$; 
\item (Unit) $\forall i$, one has $b_0 b_i=b_ib_0=b_i$; 
\item (Dual) $\forall i, \exists! i^*\in\{0,\dots,r-1\}, \forall j$ one has $N^{0}_{ij}=N^0_{ji}=\delta_{i^*j}$, where $\delta_{lm}$ is the Kronecker symbol; 
\item (Frobenius reciprocity) $\forall i,j,k$, one has $N^k_{ij}=N^j_{i^*k}=N^i_{kj^*}$. 
\end{itemize}
Put $N_i:=(N_{ij}^k)_{j,k}$. Then $N_i^{\mathsf T}=N_{i^*}$. 
We say that $R$ is \textit{self-dual} if $i^*=i$ for all $i$. 
The \textit{Frobenius--Perron dimension} is the unique ring homomorphism $\FPdim:R\to\mathbb R$ that is positive on the distinguished basis. 
The spectral radius of $N_i$ is $d_i:=\FPdim(b_i)$ \cite[Theorem~8.2 and Lemma~8.3]{ENO}.
Write
\[
 \FPdim(R)=\sum_i d_i^2 \qquad \text{and} \qquad \mu(R)=\max_{i,j,k}N_{ij}^k.
\]

A fusion ring of rank at most five is commutative by \cite[Proposition~9.2]{IntegralRings}. 
Thus, what follows is valid for $r=4$ or $r=5$. 

Let $R$ be a commutative fusion ring of rank $r$.
It then has $r$ characters, denoted by $\chi_0,\dots, \chi_{r-1}$, and one of them is the homomorphism $\FPdim$, say $\chi_0$ \cite[Proposition 3.3.6 (3)]{EGNO}.
Given $j$, the \textit{formal codegree} of $\chi_j$ is the real number $f_j=\sum_{i=0}^{r-1} \vert \chi_j(b_i)\vert^2$. 
Since $\chi_j(b_i)$ is an eigenvalue of $N_i$ for every $i$, we infer that $f_j$ is an algebraic integer and that $\vert \chi_j(b_i)\vert \leq d_i$; whence $f_j \leq \FPdim(R)= f_0$. 
Also, $\chi_j$ is a real character if $R$ is self-dual, since $N_i$ is then a real symmetric matrix for every $i$. 

The normalized character vectors
$(\chi_j(b_i))_i/\sqrt{f_j}$ form an orthonormal basis, giving
\begin{equation}\label{eq:orthogonality}
 \sum_{j=0}^{r-1}\frac{\chi_j(b_i)\overline{\chi_j(b_k)}}{f_j}
 =\delta_{ik} \qquad  \text{and} \qquad \sum_{j=0}^{r-1} \frac{1}{f_j}=1.
\end{equation}
See \cite[Theorem 2.4]{LPR2} for the first equality. 
The second is obtained from the first one by taking $i=k=0$. In particular, $f_j>1$ if $r\geq2$. 

\begin{lemma}\label{lem:commutative}
Let $R$ be a commutative fusion ring of rank $r$. Its fusion matrices are
normal and the eigenvalues of its \textnormal{Casimir matrix} $\sum_iN_iN_i^{\mathsf T}$ are the formal codegrees $f_0,\dots, f_{r-1}$.
In addition, $\mu(R)^2\leq\FPdim(R)$.
\end{lemma}
\begin{proof}
The Casimir spectrum is given by
\cite[Example~2.4 and Lemma~2.6]{OstrikCodegrees}.
The normality of fusion matrices arises from their pairwise commutativity and from $N_i^{\mathsf T}=N_{i^*}$. 
Hence, the Euclidean operator norm of $N_i$ equals the spectral radius of $N_i$, namely $d_i$, and so $N_{ij}^k\leq d_i\leq\sqrt{\FPdim(R)}$.
\end{proof}

\subsection{Generalities on fusion categories} \label{sub-sec2.2}
For a fusion category $\mathcal C$, let $\mathcal O(\mathcal C)$ be a
set of representatives of its simple isomorphism classes and write
$V_0=\one,\ldots,V_{r-1}$ for these representatives.
Its Grothendieck ring $K_0(\mathcal C)$ is a fusion ring with distinguished basis
$\{1, [V_1],\dots, [V_{r-1}]\}$ and with multiplication induced by $\otimes$
\cite[Proposition~4.9.1]{EGNO}. We set
$\FPdim(\mathcal C)=\FPdim(K_0(\mathcal C))$.
Its \textit{categorical dimension} is
\[
 \dim(\mathcal C)=\sum_{V\in\mathcal O(\mathcal C)}|V|^2,
\]
where $|V|^2$ is the squared norm of $V$ \cite[Definition~7.21.3]{EGNO}.
We say that $\mathcal{C}$ is \emph{pseudo-unitary} if
$\dim(\mathcal C)=\FPdim(\mathcal C)$.
It is known that $\mathcal{C}$ is pseudo-unitary if $\FPdim(\mathcal C)\in\Z$ \cite[Proposition~9.6.5]{EGNO}.

Applying any field automorphism $\sigma$ of $\C$ to the
structure morphisms gives the Galois-conjugate fusion category
$\mathcal C^\sigma$. The fusion coefficients and the set of simple objects remain
unchanged, so $K_0(\mathcal C^\sigma)\cong K_0(\mathcal C)$ and
$\FPdim(\mathcal C^\sigma)=\FPdim(\mathcal C)$.
One also has $\dim(\mathcal C^\sigma)=\sigma(\dim(\mathcal C))$
\cite[\S2.1]{OstrikGlobal}.

A \textit{pivotal structure} on $\mathcal C$ is a monoidal natural isomorphism
$j:\operatorname{Id}_{\mathcal C}\Rightarrow(-)^{**}$
\cite[Definition~4.7.7]{EGNO}. It defines a dimension character $\chi$ on $K_0(\mathcal{C})$ by $\dim_{\mathcal C}(V):=\chi([V])$, which is nonzero on every simple object $V$ \cite[Propositions~4.7.12 and~4.8.4]{EGNO}. 
Moreover, \[
 \dim(\mathcal C)=
 \sum_{V\in\mathcal O(\mathcal C)}
 \dim_{\mathcal C}(V)\dim_{\mathcal C}(V^*).
\]
The pivotal structure is said to be \textit{spherical} if $\dim_{\mathcal C}(V)=\dim_{\mathcal C}(V^*)$ for every simple object $V$
\cite[Definition~4.7.14]{EGNO}. 
For a commutative Grothendieck ring, this
is equivalent to these dimensions being real. 
A pseudo-unitary fusion category admits a canonical spherical structure
with simple dimensions equal to $1,\FPdim([V_1]), \dots, \FPdim([V_{r-1}])$ \cite[Proposition~9.5.1]{EGNO}.

\subsection{Reductions} 
Theorem \ref{thm:uniform} establishes that any fusion ring $R$ of rank four admitting a pivotal categorification must satisfy $\FPdim(R) < 3600$.
To obtain such a bound, we need to reduce the problem to a setting with more favorable hypotheses. 
This is the purpose of this subsection. 
For example, for rank-four fusion rings, pivotal categorifiability and spherical categorifiability coincide.

\begin{lemma}\label{lem:spherical}
Every pivotal fusion category of rank four admits a spherical structure.
\end{lemma}
\begin{proof}
A pivotal fusion category admitting no spherical structure has at
least five isomorphism classes of simple objects \cite[\S7, Lemma~4]{HaggeHong}.
\end{proof}

We now explain that if $\dim(\mathcal{C})$ and $\FPdim(\mathcal{C})$ have the same minimal polynomial over $\Q$, then we may assume without loss of generality that $\mathcal{C}$ is pseudo-unitary. 

\begin{lemma}\label{lem:galois-reduction}
Let $\mathcal C$ be a fusion category.
If $\dim(\mathcal C)$ and $\FPdim(\mathcal C)$ are Galois conjugate,
then $\mathcal C$ has a pseudo-unitary Galois conjugate.
\end{lemma}
\begin{proof}
Choose a field automorphism $\sigma$ of $\C$ with $\sigma(\dim(\mathcal C))=\FPdim(\mathcal C)$.
The preceding identities give
$\dim(\mathcal C^\sigma)=\sigma(\dim(\mathcal C))=\FPdim(\mathcal C)=\FPdim(\mathcal C^\sigma)$.
\end{proof}

\subsection{Restrictions} \label{sub-sec2.4}
As we will see in Section \ref{sec:established}, there are $8977$ fusion rings of rank four whose global Frobenius--Perron dimension is less than $3600$. 
We can now check one by one whether these fusion rings admit a pivotal categorification. 
This subsection reviews exclusion criteria that reduce the list of potential candidates to $75$.
The indicator inequalities of the next section will rule out another $60$.

\begin{theorem} \label{prop:subring-Fibo} \textnormal{(Ostrik, \cite[Main Theorem, p.~178]{OstrikRank2})}
A fusion ring with basis $\{1,x\}$ and fusion rule $x^2=1+ax$ is
categorifiable only for $a=0,1$. 
\end{theorem}

A nonzero algebraic integer $\alpha$ is called a \textit{unit} if $\alpha^{-1}$ is an algebraic integer and a \textit{$d$-number} if $\beta/\alpha$ is a unit for all Galois conjugates $\beta$ of $\alpha$. 
Letting $n$ be the degree of $\alpha$ over $\Q$, the element $\alpha$ is a $d$-number if and only if $\alpha^n=k\varepsilon$ for some nonzero integer $k$ and some unit $\varepsilon$ \cite[Lemma~2.7(iv)]{OstrikCodegrees}. 

\begin{theorem} \label{prop:d-number} \textnormal{(Ostrik, \cite[Theorem~1.2]{OstrikCodegrees})}
Every formal codegree of a categorifiable fusion ring is a $d$-number. 
\end{theorem} 

We also use three published rank-four classifications. In each nonpointed
family below, the basic Frobenius--Perron dimensions take exactly the
values $1,d$, with $d>1$. Any categorification is therefore Galois
conjugate to a pseudo-unitary one \cite[Lemma~3.1]{Schopieray}.
The following are necessary conditions for categorifiability:
\begin{enumerate}
\item For the near-group ring $\Near(C_3,m)$, with basis $C_3\sqcup\{x\}$
and rules $gx=xg=x$, $x^2=\sum_{g\in C_3}g+mx$, one has
$m\in\{0,2,3,6\}$ \cite[Theorem~1.1]{Larson}.
\item For the ring with basis $\{1,g,x,x^*\}$ and rules
$g^2=1$, $gx=x^*$, $x^2=g+m(x+x^*)$, one has $m\leq2$
\cite[Theorem~1.1]{Larson}.
\item For the self-dual ring with basis $\{1,g,x,y\}$ and rules
$g^2=1$, $gx=y$, $x^2=1+ax+by$, one has
$(a,b)\in\{(0,0),(0,1),(1,0),(1,1),(2,2)\}$
\cite[Theorem~2.4]{EIP}.
\end{enumerate}

Up to fusion ring isomorphism, Theorem \ref{thm:main} classifies all fusion rings of rank four admitting a pivotal categorification.
Any such categorification must be spherical by Lemma \ref{lem:spherical}. 
Hence, for our purposes, it is particularly relevant to study the properties of a spherical fusion category $\mathcal{C}$ whose Grothendieck ring is commutative, as we do in the remainder of this subsection. 

Put $c_i=\dim_{\mathcal C}(V_i)$ for every $i\in\{0,\dots,r-1\}$. 
Note that $c_i=c_{i^*}$ for all $i$ and that $\dim(\mathcal C)=\sum_i c_i^2$. 
The $c_i$ are nonzero totally real algebraic integers \cite[\S2.1]{OstrikRank3}.
Ostrik's spherical inequality yields
\begin{equation}\label{eq:spherical-mass}
\sum_j \frac{1}{f_j^2} \leq\frac12+\frac1{2\sigma(\dim(\mathcal C))}
\end{equation}
for all field automorphisms $\sigma$ of $\C$ \cite[Theorem~3.1.1 and Remark~3.1.2]{OstrikGlobal}.  
Also, $\dim(\mathcal C)/f$ must be a cyclotomic integer for every formal codegree $f$ of $K_0(\mathcal{C})$ \cite[Theorem~2.13]{OstrikRank3}.
In particular, the norm of $\dim(\mathcal C)/f$ is an integer and the field generated by $\dim(\mathcal C)/f$ is abelian over $\Q$. 
For a ratio of formal codegrees, these conditions also suffice for cyclotomic integrality, as the following lemma shows. 

\begin{lemma}\label{lem:codegree-ratio}
Let $u,v$ be $d$-numbers and put $x=u/v$. Then $x$ is a cyclotomic
integer if and only if $N_{\Q(x)/\Q}(x)\in\Z$ and $\Q(x)/\Q$ is abelian.
\end{lemma}
\begin{proof}
The forward implication is clear. Conversely, extend each embedding of
$\Q(x)$ to an automorphism $\sigma$ of $\C$. Since $u,v$ are $d$-numbers,
$\sigma(x)/x=(\sigma(u)/u)(v/\sigma(v))$ is an algebraic-integer unit.
Thus $N_{\Q(x)/\Q}(x)$ is $x^{[\Q(x):\Q]}$ times a unit. Its integrality
implies that $x$ is an algebraic integer. The abelianity assumption and
the Kronecker--Weber theorem place $x$ in a cyclotomic field.
\end{proof}

The abelianity condition can be checked once for the entire ring.
Indeed, for any formal codegree $H$, equation~\eqref{eq:orthogonality}
gives $H=\sum_jH/f_j$. Hence all fields $\Q(H/f_j)$ are abelian
if and only if the splitting field of $\prod_j(t-f_j)$ is abelian.
This polynomial is the characteristic polynomial of the Casimir matrix
and has degree four in our computation.

\section{Central induction and indicator inequalities}\label{sec:central}
Let $\mathcal C$ be a pseudo-unitary fusion category, equipped with its
canonical spherical structure, and put $D=\FPdim(\mathcal C)$.
Let $F:\mathcal Z(\mathcal C)\to\mathcal C$ be the forgetful functor
and $I$ its right adjoint. For a simple object $A$ of the center, write
$\theta_A$ for its twist and $[I(V):A]$ for its multiplicity in $I(V)$.
We use the following identities:
\[
 \sum_A[I(V):A]\FPdim(A)=D\FPdim(V),\qquad
 \sum_A[I(V):A]\FPdim(A)\theta_A^n=D\nu_n(V)\quad(n\geq1).
\]
The first follows from
$FI(V)\cong\bigoplus_{X\in\mathcal O(\mathcal C)}X\otimes V\otimes X^*$
\cite[Proposition~2.4]{OstrikRank3}; the second is the indicator formula
\cite[Theorem~4.1]{NgSchauenburg}. The twists are roots of unity
\cite[Theorem~5.5]{NgSchauenburg}, and all simple summands of $I(1)$
have twist one \cite[Theorem~2.5]{OstrikRank3}.

The \emph{low part} of $I(V)$ consists of its simple summands that also
occur in $I(1)$, with their multiplicities in $I(V)$:
\[
 I_{\mathrm{low}}(V)
 =\bigoplus_{A\in\mathcal O(I(1))}A^{\oplus[I(V):A]}.
\]
Here $\mathcal O(I(1))$ denotes the simple support of $I(1)$.
Write $d_V=\FPdim(V)$ and $W_V=\FPdim(I_{\mathrm{low}}(V))$.
For a root of unity $\zeta$, let $I(V)_\zeta$ be the sum of the simple
summands of $I(V)$ with twist $\zeta$, again with their multiplicities.
Thus $I_{\mathrm{low}}(V)$ is a direct summand of $I(V)_1$.

\begin{proposition}[Polynomial indicator bounds]\label{lem:indicator-polynomials}
Let $V\in\mathcal C$ and let
\[
 P(z)=c_0+\sum_{k=1}^n(c_kz^k+\overline{c_k}z^{-k}),
 \qquad c_0\in\mathbb R,\quad c_k\in\C,
\]
be nonnegative on the unit circle. Then
\[
 P(1)W_V\leq D\left(c_0d_V+
                    2\operatorname{Re}\sum_{k=1}^n c_k\nu_k(V)\right).
\]
For every root of unity $\zeta$, one also has
\[
 P(1)\FPdim(I(V)_\zeta)
 \leq D\left(c_0d_V+
       2\operatorname{Re}\sum_{k=1}^n c_k\zeta^{-k}\nu_k(V)\right).
\]
\end{proposition}
\begin{proof}
The indicator formula and $|\theta_A|=1$ give
\[
 \sum_A[I(V):A]\FPdim(A)P(\theta_A)
 =D\left(c_0d_V+2\operatorname{Re}\sum_{k=1}^n c_k\nu_k(V)\right).
\]
Every term on the left is nonnegative. The summands from $I(1)$ all
have twist one, so keeping only those terms gives the first inequality.
For the second, replace $P(z)$ by $P(\zeta^{-1}z)$ and keep the terms
with $\theta_A=\zeta$.
\end{proof}

We obtain the polynomials used below by taking squared moduli.
For $Q(z)=\sum_{j=0}^n a_jz^j$ and $|z|=1$,
\[
 |Q(z)|^2=c_0+\sum_{k=1}^n(c_kz^k+\overline{c_k}z^{-k}),\qquad
 c_0=\sum_{j=0}^n|a_j|^2,\quad
 c_k=\sum_{j=0}^{n-k}a_{j+k}\overline{a_j}.
\]
Here $\overline z=z^{-1}$; this Laurent expansion is an identity on
the unit circle.

\begin{remark}[Vanishing indicators]\label{rem:vanishing-indicators}
If $\nu_1(V)=\cdots=\nu_m(V)=0$ for an integer $m\geq1$, apply
Proposition~\ref{lem:indicator-polynomials} to
\[
 |1+z+\cdots+z^m|^2
 =(m+1)+\sum_{k=1}^m(m+1-k)(z^k+z^{-k})
\]
to obtain
\[
 (m+1)W_V\leq Dd_V,\qquad
 (m+1)\FPdim(I(V)_\zeta)\leq Dd_V.
\]
For $m=1$, the low-part bound is the estimate in
\cite[proof of Theorem~2.21]{OstrikRank3}; the extension to twist
sectors was suggested by Andrew Schopieray. For a direct sum of
non-self-dual simples, one can take $m=2$, by additivity of indicators
\cite[Corollary~7.8]{NgSchauenburgHigher}.
\end{remark}

\begin{corollary}\label{cor:second-indicator}
For a nonunit simple object $X$,
\[
 9W_X\leq D\bigl(3d_X+2\nu_2(X)\bigr).
\]
In particular,
\[
 9W_X\leq D(3d_X+2)\quad(X\cong X^*),\qquad
 3W_X\leq Dd_X\quad(X\not\cong X^*).
\]
\end{corollary}
\begin{proof}
Apply Proposition~\ref{lem:indicator-polynomials} to
$|1+z+z^2|^2=3+2(z+z^{-1})+(z^2+z^{-2})$.
Here $\nu_1(X)=0$, while $\nu_2(X)$ is $\pm1$ for a self-dual simple
and zero otherwise \cite[Definition~3.1 and Theorem~5.1]{NgSchauenburgHigher}.
\end{proof}

\begin{example}\label{ex:S(2)}
Corollary~\ref{cor:second-indicator} rules out the ring $S(2)$ of
\cite[\S4, Theorem~4.1]{EIP}, already listed by Larson
\cite[Theorem~1.1(6)]{Larson} as $K_2(2)$. It has basis
$\{1,\alpha,\rho,\rho^*\}$, with
$\alpha^2=1$, $\rho^*=\alpha\rho$, and
$\rho^2=\alpha+2\rho+2\rho^*$.
Put $d=2+\sqrt5$ and $D=4+8d$; its simple Frobenius--Perron
dimensions are $1,1,d,d$. Any categorification has a pseudo-unitary
Galois conjugate \cite[Lemma~3.1]{Schopieray}. Larson's induction
formulas \cite[Proposition~5.2]{Larson}, with $c=2$, then give
\[
 W_\rho=2(1+4d)+2(1+2d)=4+12d.
\]
Since $\rho$ is not self-dual, the corollary would imply
\[
 12+36d=3W_\rho\leq Dd=8+36d,
\]
a contradiction. Thus $S(2)$ has no complex fusion categorification.
\end{example}

\begin{corollary}\label{cor:third-indicator}
For a nonunit simple object $X$, put
$a=\dim_{\C}\operatorname{Hom}(1,X^{\otimes3})=N_{XX}^{X^*}$. Then
\[
 36W_X\leq D(10d_X+8+2a).
\]
\end{corollary}
\begin{proof}
Apply Proposition~\ref{lem:indicator-polynomials} to
\[
 |(1+z)(1+z+z^2)|^2
 =10+8(z+z^{-1})+4(z^2+z^{-2})+(z^3+z^{-3}).
\]
Use $\nu_1(X)=0$, $\nu_2(X)\leq1$ and $|\nu_3(X)|\leq a$.
The last bound holds because $\nu_n(X)$ is the trace of a rotation on
$\operatorname{Hom}(1,X^{\otimes n})$ whose $n$th power is the identity
\cite[Definition~3.1 and Theorem~5.1]{NgSchauenburgHigher}.
\end{proof}

\begin{corollary}[Fourth-indicator bound]\label{prop:fourth-indicator}
For a nonunit simple object $X$, put
\[
 s=\dim_{\C}\operatorname{Hom}(1,X^{\otimes2})=\mathbf1_{X\cong X^*},
 \qquad a=\dim_{\C}\operatorname{Hom}(1,X^{\otimes3})=N_{XX}^{X^*}.
\]
Then
\[
 32W_X\leq D\left(7d_X+
       \left\lfloor\frac{10d_X+71s+30a}{9}\right\rfloor\right).
\]
\end{corollary}
\begin{proof}
The rotation eigenvalues give $\nu_2(X)\in\Z$,
$2\operatorname{Re}\nu_3(X)\in\Z$ and
$\operatorname{Re}\nu_4(X)\in\Z$. Hence
\[
 J:=8\nu_2(X)+4\operatorname{Re}\nu_3(X)
                     +\operatorname{Re}\nu_4(X)\in\Z.
\]
Expanding
\[
 P_+(z)=|(1+z)^2(1+z^2)|^2,\qquad
 P_-(z)=|3+z-z^3-3z^4|^2
\]
and applying Proposition~\ref{lem:indicator-polynomials} gives
\[
 32W_X\leq D(7d_X+J),\qquad
 \nu_2(X)+6\operatorname{Re}\nu_3(X)+9\operatorname{Re}\nu_4(X)
 \leq10d_X.
\]
Since $\nu_2(X)\leq s$ and $\operatorname{Re}\nu_3(X)\leq a$,
\[
 9J\leq10d_X+71\nu_2(X)+30\operatorname{Re}\nu_3(X)
     \leq10d_X+71s+30a.
\]
Taking the integer part proves the claim.
\end{proof}

To apply these inequalities using fusion rules, suppose now that
$K_0(\mathcal C)$ is commutative of rank $r$, with simples
$X_0=1,X_1,\ldots,X_{r-1}$. Put $d_i=\FPdim(X_i)$,
$d=(d_i)_{i=0}^{r-1}$ and $B=\sum_iN_iN_i^{\mathsf T}$.
The \emph{induction matrix} has entries
$M_{ai}=[I(X_i):Z_a]=[F(Z_a):X_i]$, where $Z_a$ ranges over the
simple objects of the center. Let $L$ consist of the rows for the
simple summands of $I(1)$. There are $r$ such summands, each occurring
once, with dimensions $D/f_j$ \cite[Theorem~2.13]{OstrikRank3}.
One row is $e_0^{\mathsf T}=(1,0,\ldots,0)$, corresponding to the
central unit. Adjunction gives
\[
 M^{\mathsf T}M=B,\qquad L^{\mathsf T}\mathbf1=Be_0,\qquad
 Q:=B-L^{\mathsf T}L,\qquad W_i:=(L^{\mathsf T}Ld)_i.
\]
Thus $W_i=W_{X_i}$ and $Q$ is the Gram matrix of the rows outside $L$.
Every row $v^{\mathsf T}$ of $L$, assigned to a formal codegree $f$, satisfies
\[
 v\in\Z_{\geq0}^r,\qquad v_0=1,\qquad d^{\mathsf T}v=D/f,
 \qquad v_i\leq\min\{B_{0i},\lfloor\sqrt{B_{ii}}\rfloor\}.
\]
The bounds follow from the first column and diagonal of
$M^{\mathsf T}M=B$. They give a finite list of possible $L$.
For each choice, the preceding corollaries can be tested directly.
Rows assigned to equal formal codegrees may be permuted without changing
$L^{\mathsf T}L$, so these permutations need not be tested separately.

\begin{example}\label{ex:R2}
Corollary~\ref{prop:fourth-indicator} rules out the self-dual fusion ring
$R_2$ with basis $\{1,X,Y,Z\}$ and products
\[
\begin{array}{lll}
 X^2=1+X+Y+2Z,&XY=X+2Y+2Z,&XZ=2X+2Y+2Z,\\
 Y^2=1+2X+2Y+2Z,&YZ=2X+2Y+3Z,&Z^2=1+2X+3Y+3Z.
\end{array}
\]
Direct multiplication checks the fusion-ring axioms. Put
$\delta=3+\sqrt{10}$, $\delta'=3-\sqrt{10}$ and $D=6+18\delta$.
The four characters are
\[
 (1,\delta-1,\delta,\delta+1),\quad
 (1,\delta'-1,\delta',\delta'+1),\quad
 (1,0,1,-1),\quad(1,1,-1,0).
\]
The first is the Frobenius--Perron character. Any pivotal structure is
spherical, since the simples are self-dual. The rational characters
vanish on a simple, so Lemma~\ref{lem:galois-reduction} reduces any
pivotal categorification to a pseudo-unitary one.
Its formal codegrees are $D,6+18\delta',3,3$.
The nonunit summands of $I(1)$ therefore have dimensions
$1+6\delta,2+6\delta,2+6\delta$, and the preceding row conditions give
\[
 (1,r,6-2r,r),\qquad (1,p,5-2p,p+1),\qquad (1,q,5-2q,q+1).
\]
Since $L^{\mathsf T}\mathbf1=Be_0$, we have $r+p+q=5$;
nonnegativity gives $r\leq3$. Hence
\[
 W_X=r(1+6\delta)+(p+q)(2+6\delta)
     =10-r+30\delta\geq7+30\delta.
\]
But $X^2=1+X+Y+2Z$ gives $d_X=\delta-1$ and $s=a=1$ in
Corollary~\ref{prop:fourth-indicator}. Since $6<\delta<31/5$,
$\lfloor(10\delta+91)/9\rfloor=16$. The corollary would give
\[
 224+960\delta\leq32W_X
 \leq D\bigl(7(\delta-1)+16\bigr)=180+960\delta,
\]
where $\delta^2=6\delta+1$. This is a contradiction.
\end{example}

One additional inequality is useful when a simple object is an invertible
involution.
\begin{proposition}\label{prop:low-test}
If $g\ne1$ is an invertible involution satisfying
$Q_{gg}=2$, $D/2-W_g=1$ and $W_g>0$, then
\[
 2W_i+DQ_{gi}\leq D(d_i+1)\qquad(i\ne0,g).
\]
\end{proposition}
\begin{proof}
Since $Q_{gg}=2$, exactly two simple summands of $I(g)$ lie outside
its low part, each with multiplicity one. The sum of all summand
dimensions is $D$, while the second twist trace is $D\nu_2(g)$.
Here $|\nu_2(g)|=1$, so equality in the triangle inequality forces
$\theta_A^2=\nu_2(g)$ for each summand $A$ of $I(g)$.
As $W_g>0$, one such twist is one. Thus $\nu_2(g)=1$ and all twists
are $\pm1$.

The first twist trace is zero, so the dimensions at each of the two
twists sum to $D/2$. Outside the low part, the dimensions are therefore
$D/2-W_g=1$ at twist $1$ and $D/2$ at twist $-1$.
The dimension-one summand restricts to $g$; the other has induction row
$(Q_{gi}-\delta_{gi})_i$.
For $i\ne0,g$, these summands and the low part consequently contribute
at least $W_i+(D/2)Q_{gi}$ to the dimension at twists $\pm1$ in $I(X_i)$.
The polynomial $1+\operatorname{Re}(z^2)$ is nonnegative on the circle
and equals two at both $1$ and $-1$. The identity in the proof of
Proposition~\ref{lem:indicator-polynomials} therefore gives
\[
 2W_i+DQ_{gi}\leq D(d_i+\nu_2(X_i))\leq D(d_i+1).
\]
\end{proof}

\begin{samepage}
\begin{lemma}\label{lem:short-moments}
Suppose $K_0(\mathcal C)$ is commutative of rank $r\geq2$, and put
\[
 S=\sum_{i>0}d_i,\qquad A=\sum_{i>0}d_iN_{ii}^{i^*},\qquad
 w=\sum_{j=0}^{r-1}f_j^{-2}.
\]
Then
\begin{equation}\label{eq:short-second-moment}
 9(D^2w-D)\leq D\bigl(3(D-1)+2S\bigr),
\end{equation}
\[
 36(D^2w-D)\leq D\bigl(10(D-1)+8S+2A\bigr).
\]
In particular, if $r=4$, then
\begin{equation}\label{eq:short-diagonal-lower}
 A\geq D-25-4S+\frac{24}{D}.
\end{equation}
\end{lemma}
\end{samepage}
\begin{proof}
Each low row has unit entry one and dimension $D/f_j$. Hence
\[
 \sum_{i>0}d_iW_i
 =\sum_j\frac D{f_j}\left(\frac D{f_j}-1\right)=D^2w-D,
\]
using $\sum_jf_j^{-1}=1$. Multiply the inequalities of
Corollaries~\ref{cor:second-indicator} and~\ref{cor:third-indicator}
by $d_i$ and sum, using $\nu_2(X_i)\leq1$, to obtain the first two claims.
The second gives $A+4S\geq18Dw-5D-13$.
Since $f_0=D$, Cauchy--Schwarz gives
\[
 w\geq\frac1{D^2}+\frac1{r-1}\left(1-\frac1D\right)^2.
\]
For $r=4$, substitution gives the last inequality.
\end{proof}

\section{A uniform dimension bound}\label{sec:uniform}
All fusion categories in this section are over $\C$ and have rank four,
unless another rank is stated explicitly. The estimates of
Section~\ref{sec:central} yield a bound independent of
any census. This bound makes the subsequent classification a finite
problem; the sharper bounds in Remark~\ref{rem:sharp} follow only
after that classification.
\begin{theorem}\label{thm:uniform}
Every rank-four fusion ring $R$ admitting a pivotal categorification satisfies
\[
 \FPdim(R)<3600,\qquad \mu(R)\leq59.
\]
\end{theorem}
\subsection{Reduction to the self-dual pseudo-unitary case}\label{sec:pivotal-reduction}

\begin{lemma}\label{lem:exceptional}
Let $\mathcal C$ be a self-dual pivotal fusion category of rank four. Either it has a pseudo-unitary Galois conjugate, or $\FPdim(\mathcal C)<576$.
\end{lemma}
\begin{proof}
Write $R=K_0(\mathcal C)$ and $D=\FPdim(\mathcal C)$.
Self-duality makes the pivotal structure spherical. If $D$ is rational,
it is an integer and the fusion category is pseudo-unitary by \cite[Proposition~8.24]{ENO}. For the rest of the proof, assume that no Galois conjugate is pseudo-unitary. Then $D$ is irrational, and the spherical dimension character is outside the Frobenius--Perron orbit.

The field generated by spherical dimensions contains every formal codegree, including $D$, by \cite[Corollary~2.15]{OstrikRank3}. Here the \emph{Frobenius--Perron field} means
$K=\Q(\FPdim(X):X\in\mathcal O(\mathcal C))$, the field of
values of the Frobenius--Perron character; it need not equal $\Q(D)$
a priori. Its degree is the size of that character orbit, hence lies
between two and four. Degree four leaves no other character orbit, and in degree three the other orbit is rational, so neither case is possible. If $[K:\Q]=2$, then $\Q(D)=K$ because $D$ is irrational.
The other character field contains $K$, and only two dimensions remain
in $R\otimes\Q$; it must therefore be another copy of $K$. Thus
\begin{equation}\label{eq:repeated-quadratic}
 R\otimes\Q\cong K\times K
\end{equation}
for a real quadratic field $K$, and the spherical dimension character belongs to the other factor.

Denote the four formal codegrees by $D,D',H,H'$, with primes indicating conjugation in $K$. We allow $H=H'$ and list one formal codegree for each character. After Galois conjugation, arrange that the spherical global dimension is $H\geq H'$. Each formal codegree is positive and greater than one.
By \cite[Propositions~8.21 and~8.22]{ENO}, $H\leq D$ and
$H/D$ is an algebraic integer. Equality would make the fusion
category pseudo-unitary, so $H<D$ in the exceptional case. Put
\[
 n=DD',\qquad k=\frac{HH'}{DD'}\in\Z_{\geq1},\qquad
 p=\frac1D+\frac1{D'},\qquad
 w=\frac1{D^2}+\frac1{(D')^2}+\frac1{H^2}+\frac1{(H')^2}.
\]
Here $k$ is the field norm of $H/D$, including when that element is rational. Since $D'>1$ and $H\geq H'$, we have
\begin{equation}\label{eq:exceptional-size}
 n>D,\qquad H\geq\sqrt{kn}>\sqrt{kD}.
\end{equation}
Equation \eqref{eq:orthogonality} gives $D^{-1}+(D')^{-1}+H^{-1}+(H')^{-1}=1$, and so
\begin{equation}\label{eq:exceptional-w}
 w-\frac12=2\left(p-\frac12\right)^2-\frac{2(1+1/k)}n.
\end{equation}

Ostrik's spherical inequality \eqref{eq:spherical-mass},
applied with global dimension $H$, gives
\[
 w\leq\frac12+\frac1{2H}.
\]

Set $r=D/H>1$. By \eqref{eq:exceptional-size}, $r<\sqrt{D/k}\leq\sqrt D$. Moreover, $kr>2$. This is immediate when $k\geq2$. Both $H/D$ and its conjugate $H'/D'$ are positive. When $k=1$,
$H/D$ is therefore a totally positive algebraic-integer unit of norm one. Thus $r$ is also such a unit, with conjugate $r^{-1}$; the integer $r+r^{-1}>2$ is at least three, and $r\geq(3+\sqrt5)/2>2$.

Since $H'=krD'$, the same identity gives
\[
 1=\frac{1+r}{D}+\frac{1+1/(kr)}{D'},\qquad
 p=\frac{kr}{kr+1}-\frac{kr^2-1}{(kr+1)D}
 >\frac23-\frac rD>\frac23-\frac1{\sqrt D}.
\]
If $D\geq576$, then $p>5/8$ and $H>24$. Equation~\eqref{eq:exceptional-w} consequently gives
\[
 w-\frac12>\frac1{32}-\frac4D
 \geq\frac7{288}>\frac1{48}>\frac1{2H},
\]
contradicting \eqref{eq:spherical-mass}. Thus $D<576$, as required.
\end{proof}

\begin{lemma}\label{lem:nonselfdual}
A pivotal fusion category of rank four with a non-self-dual simple has Frobenius--Perron dimension less than $45$.
\end{lemma}
\begin{proof}
A pointed fusion category has Frobenius--Perron dimension four, so assume
that the fusion category is not pointed. Write the basis as $1,X,Y,Z$, where $X^*=Z$ and $Y^*=Y$.

By Lemma~\ref{lem:spherical}, we may choose a spherical structure.

There are exactly two real-valued characters of this fusion ring. Indeed, by the adjoint relation $N_{U^*}=N_U^{\mathsf T}$ and simultaneous unitary diagonalization, $\chi(U^*)=\overline{\chi(U)}$ for every basis element $U$. Extend basis duality $\C$-linearly. This involution has trace two on the distinguished basis; on the primitive idempotents it is a permutation, so exactly two characters are fixed, precisely the real-valued ones. If the Frobenius--Perron character is irrational, its Galois orbit has at least two members, all fixed by duality because
$\FPdim(U)=\FPdim(U^*)$, and hence consists of these two characters. The spherical dimension character is real-valued, so it belongs to that orbit. Lemma~\ref{lem:galois-reduction} gives a pseudo-unitary Galois conjugate. If the Frobenius--Perron character is rational, the global Frobenius--Perron dimension is an integer and pseudo-unitarity follows directly from \cite[Proposition~8.24]{ENO}.

Larson's candidate classification \cite[Theorem~1.1]{Larson} bounds
the global Frobenius--Perron dimension of a non-self-dual pseudo-unitary
fusion category of rank four by $24+12\sqrt3<45$.
The Galois conjugate has the same fusion ring, so the same bound holds
for $\mathcal C$.
\end{proof}

\subsection{Concentration in the self-dual case}\label{sec:concentration}
Let $\mathcal C$ be a self-dual pseudo-unitary fusion category of rank four,
with simples $1,X_1,X_2,X_3$. Put
\[
 d_i=\FPdim(X_i),\quad a_i=N_{ii}^{i},\quad
 S=\sum_{i=1}^3d_i,\quad A=\sum_{i=1}^3d_i a_i,\quad
 w=\sum_{j=0}^3f_j^{-2}.
\]
Its fusion ring is commutative. Index the characters so that
$\chi_0=\FPdim$ and $f_0=\FPdim(\mathcal C)$.
Character orthogonality gives
\begin{equation}\label{eq:short-character-identities}
 a_i=\sum_{j=0}^3\frac{\chi_j(X_i)^3}{f_j},
\end{equation}
by taking the coefficient of $1$ in $X_i^3$.

\begin{lemma}\label{lem:short-concentration}
Suppose that $\FPdim(\mathcal C)\geq3600$. Put
$M=\max_i d_i^2$ and $E=\FPdim(\mathcal C)-M$. Then
\begin{gather}
 2<f_j<4\quad(1\leq j\leq3),\qquad f_1f_2f_3<36,
       \label{eq:short-codegrees}\\
 E<\frac72\sqrt{\FPdim(\mathcal C)}.
       \label{eq:short-dominance}
\end{gather}
There is a unique simple object $X$ of maximum dimension. Its fusion
matrix is singular, and
\begin{equation}\label{eq:short-small-eigenvalues}
 |\chi_j(X)|^2<\frac{14}{\sqrt{\FPdim(\mathcal C)}}
 \leq\frac7{30}<\frac14
 \qquad(1\leq j\leq3).
\end{equation}
\end{lemma}

\begin{proof}
Put $D=\FPdim(\mathcal C)$, $t=\sqrt D\geq60$, $q_j=f_j^{-1}$,
$m=(1-D^{-1})/3$, and $V=\sum_{j=1}^3(q_j-m)^2$.
Substitution of $w=D^{-2}+3m^2+V$ into
\eqref{eq:short-second-moment} gives
\[
 V\leq\frac4{3D}+\frac{2S}{9D}-\frac4{3D^2}
 <\frac4{3t^2}+\frac7{18t}\leq\frac{37}{5400}<\frac1{144},
\]
where $S<\sqrt3t<7t/4$.
Three real numbers of sum zero each have square at most two thirds of
their squared norm. Therefore
\[
 |q_j-m|\leq\sqrt{2V/3}<\frac1{14},\qquad
 \frac14<\frac{11}{42}-\frac1{3D}<q_j
 <\frac{17}{42}<\frac12.
\]
This proves $2<f_j<4$. Writing $q_j=1/4+r_j$ with $r_j>0$ gives
\[
 q_1q_2q_3
 \geq\frac1{64}+\frac{r_1+r_2+r_3}{16}
 =\frac1{32}-\frac1{16D}>\frac1{36},
\]
and hence the product bound.

Relabel so that $d_3^2=M$. Since $f_j<4$, every non-Frobenius--Perron
character value on a nonunit simple has absolute value less than $2$.
Equations~\eqref{eq:orthogonality} and \eqref{eq:short-character-identities} imply
\begin{equation}\label{eq:short-spectral-upper}
 A\leq\frac{\sum_i d_i^4}{D}
    +2\sum_i d_i\left(1-\frac{d_i^2}{D}\right).
\end{equation}
Initially, $\sum_i d_i^4/D\leq M$ and
\eqref{eq:short-diagonal-lower} give
\begin{equation}\label{eq:short-coarse-dominance}
 E\leq25+6S<25+12t<13t.
\end{equation}
The coarse bound confines $E$ to an interval on which the following
comparison is effective; it is not yet the bound needed for the
determinant argument. For a sharper estimate, use
\[
 S\leq t+\sqrt{2E},\qquad
 \sum_i d_i^4\leq(D-E)^2+(E-1)^2,
\]
and
\[
 \sum_i d_i\left(1-\frac{d_i^2}{D}\right)
 \leq\frac{E}{t}+\sqrt{2E}.
\]
Here the largest dimension contributes $d_3E/D\leq E/t$, and the
other two contribute at most $d_1+d_2\leq\sqrt{2E}$.
Combining these estimates with \eqref{eq:short-diagonal-lower} and
\eqref{eq:short-spectral-upper} gives $g(E)\leq-(2E+23)/t^2<0$, and in particular
\begin{equation}\label{eq:short-g}
 g(E):=2E-\frac{2E^2}{t^2}-\frac{2E}{t}
              -6\sqrt{2E}-4t-25\leq0.
\end{equation}
Suppose $E\geq E_0:=7t/2$. On $[E_0,13t]$,
\[
 g'(E)=2-\frac{4E}{t^2}-\frac2t-\frac{3\sqrt2}{\sqrt E}
 >2-\frac{54}{t}-\frac6{\sqrt{7t}}>0.
\]
But
\[
 (3t-57)^2-252t=9\bigl((t-60)(t-6)+1\bigr)>0,
\]
so
\[
 g(E_0)=3t-\frac{113}{2}-6\sqrt{7t}>\frac12.
\]
This contradicts \eqref{eq:short-g} and proves
\eqref{eq:short-dominance}. In particular $E<D/2$, so $M>D/2$ and
the largest simple is unique.

Write $X=X_3$ and let $N_X$ be its fusion matrix. Put
$y_j=\chi_j(X_3)^2/f_j$ for $1\leq j\leq3$.
Orthogonality gives $y_1+y_2+y_3=E/D$. The arithmetic--geometric mean
inequality now shows
\begin{align*}
 |\det N_X|^2
 &=M(f_1f_2f_3)y_1y_2y_3\\
 &<36D\left(\frac{E}{3D}\right)^3
 <\frac{343}{6\sqrt D}\leq\frac{343}{360}<1.
\end{align*}
Thus $\det N_X=0$, since it is an integer. Finally,
$\chi_j(X_3)^2/f_j\leq E/D$ and $f_j<4$ imply
\eqref{eq:short-small-eigenvalues}.
\end{proof}

\subsection{The self-dual pseudo-unitary bound}\label{sec:pu}

\begin{proposition}\label{prop:pu}
If a pseudo-unitary fusion category has rank four and all its simple objects are self-dual, then its Frobenius--Perron dimension is less than $3600$.
\end{proposition}
\begin{proof}
Put $R=K_0(\mathcal C)$ and $D=\FPdim(\mathcal C)$, and suppose
$D\geq3600$.
Set $t=\sqrt D$, and let $X$ be the largest simple supplied by Lemma~\ref{lem:short-concentration}. Write the basis as $1,X,Y,Z$, with dimensions $1,x,y,z$, and put
\[
 K=\Q(x,y,z),\qquad r=[K:\Q],\qquad E=D-x^2.
\]

\smallskip\noindent
\emph{The rank of $N_X$ is $r$.}
The rational semisimple commutative algebra $R\otimes\Q$ is a product of number fields. Its character orbits are the embeddings of those fields, and the Frobenius--Perron orbit corresponds to $K$. If a character $\chi$ lies outside this orbit, each embedding of
$\Q(\chi(X))$ extends to its character field. Thus all conjugates of
$\chi(X)$ occur in the same non-Frobenius--Perron character orbit and
have absolute value less than one by \eqref{eq:short-small-eigenvalues}.
The norm of a nonzero algebraic integer has absolute value at least one,
so $\chi(X)=0$. On the Frobenius--Perron orbit the values are embeddings of the nonzero element $x\in K$, and are nonzero. Simultaneous diagonalization identifies these
character values, with their multiplicities, with the eigenvalues of
$N_X$. There are exactly $r$ embeddings in the Frobenius--Perron orbit,
so exactly $r$ eigenvalues are nonzero. Hence
\begin{equation}\label{eq:rank-degree}
 \operatorname{rank} N_X=r.
\end{equation}
The principal block of $N_X$ on $1,X$ is
$\left(\begin{smallmatrix}0&1\\1&a\end{smallmatrix}\right)$, so its rank is at least two. Its determinant vanishes by Lemma~\ref{lem:short-concentration}. Thus $r=2$ or $3$.

\smallskip\noindent
\emph{Rank two is impossible.}
Write
\[
 N_X=\begin{pmatrix}
 0&1&0&0\\
 1&a&b&c\\
 0&b&e&h\\
 0&c&h&k
 \end{pmatrix}.
\]
The upper left block is invertible and its Schur complement is exactly the lower right block. The Schur-complement rank formula \cite[\S0.8.5(b)]{HornJohnson}
therefore gives $e=h=k=0$. Consequently $XY=bX$ and $XZ=cX$, so $y=b$ and $z=c$ are positive integers. The second character in the Frobenius--Perron orbit has formal codegree
\[
 1+(x')^2+b^2+c^2<4.
\]
It follows that $b=c=1$. Thus $Y$ and $Z$ are distinct self-dual invertible objects
\cite[Corollary~3.3.10]{EGNO}. Their product is a fourth invertible basis element, necessarily $X$: it cannot be $1,Y$, or $Z$ by cancellation. This gives $D=4$, a contradiction.

\smallskip\noindent
\emph{Rank three is impossible.}
Now $R\otimes\Q\cong\Q\times K$ with $[K:\Q]=3$. The remaining character $\rho$ is integer-valued and has formal codegree strictly between two and four, hence equal to three. Since $\rho(X)=0$, its values on $Y,Z$ are each $\pm1$. Orthogonality with the positive character excludes equal signs. Relabeling gives
\begin{equation}\label{eq:cubic-rational}
 \rho=(1,0,1,-1),\qquad z=y+1.
\end{equation}
Applying $\rho$ to products and using Frobenius reciprocity gives
nonnegative integers $b,c,s$ such that
\begin{align}
 XY&=bX+cY+cZ, &Y^2&=1+cX+sY+sZ,\label{eq:cubic-products}\\
 YZ&=cX+sY+(s+1)Z.\notag
\end{align}
The trace of $N_Y$ is $b+2s+1$. Subtracting $\rho(Y)=1$ gives
\begin{equation}\label{eq:cubic-trace}
 \Tr_{K/\Q}(y)=b+2s.
\end{equation}

By \eqref{eq:short-dominance},
\[
 E=1+y^2+(y+1)^2<\frac72t,
 \qquad x^2=t^2-E>(t-4)^2.
\]
The dimension equation for $Y^2$ gives $cx<y^2<E/2$, whence
\[
 c<\frac{7t}{4(t-4)}<2.
\]
Thus $c\leq1$. The equation for $XY$ gives
\[
 \delta:=y-b=\frac{c(2y+1)}x\geq0,\qquad
 \delta^2\leq\frac{2E-3}{D-E}<\frac7{t-7/2}<1.
\]
For either nonidentity embedding of $K$, write $u=y'$ and $v=x'$.
By \eqref{eq:short-small-eigenvalues}, $|v|<1/2$. Applying the embedding
to the equation for $Y^2$ gives
\[
 u^2=1+cv+s(2u+1).
\]
If $u\geq-1/2$, then $c\leq1$ and $s\geq0$ give $u^2>1/2$,
hence $u>1/\sqrt2$. Its formal codegree would satisfy
\[
 f'=1+v^2+u^2+(u+1)^2>3+\sqrt2>4,
\]
contrary to \eqref{eq:short-codegrees}. Thus $y',y''<-1/2$, and
\eqref{eq:cubic-trace} gives
\[
 2s=(y-b)+y'+y''<1-1=0,
\]
a contradiction. This completes the proof.
\end{proof}

\begin{proof}[Proof of Theorem~\ref{thm:uniform}]
Let $\mathcal C$ be a pivotal categorification of $R$ and put
$D=\FPdim(\mathcal C)$. If it has a non-self-dual simple,
Lemma~\ref{lem:nonselfdual} gives $D<45$. Otherwise,
Lemma~\ref{lem:exceptional} either gives $D<576$ or supplies a
pseudo-unitary Galois conjugate with the same fusion ring.
Proposition~\ref{prop:pu} then gives $D<3600$.
The preliminary bound $\mu(R)^2\leq D$ implies $\mu(R)<60$, hence
$\mu(R)\leq59$ since fusion coefficients are integers.
\end{proof}

\section{Proof of Theorem \ref{thm:main}}\label{sec:established}

The fifteen Grothendieck rings arising from Table \ref{tab:models} are pairwise nonisomorphic: the global Frobenius--Perron dimension is an invariant of fusion rings, and the fusion categories $\mathrm{Vec}_{C_2\times C_2}$ and $\mathrm{Vec}_{C_4}$, as well as $PSU(2)_6$ and the principal even part of $S'$, do not have the same fusion rules. 
Thanks to Lemma \ref{lem:spherical}, it is therefore enough to show that there are, up to isomorphism, at most fifteen fusion rings of rank four admitting a spherical categorification. 

By Theorem \ref{thm:uniform}, such a ring must have global Frobenius--Perron dimension less than $3600$. 
We show below that there are $8977$ fusion rings of rank four whose global Frobenius--Perron dimension is less than $3600$.
To prove Theorem \ref{thm:main}, we therefore need to rule out all but fifteen of these fusion rings. 
This is done using the exclusion criteria stated in Sections \ref{sec:spherical} and \ref{sec:central}. 

The enumeration uses the following elementary criterion for associativity. 

\begin{lemma}\label{lem:commuting-multiplication}
Let $A$ be a commutative ring and let $M$ be a $A$-module with a
$A$-bilinear multiplication and a two-sided identity $1$. For $x\in M$,
write $L_x$ for the operator mapping $y\in M$ to $xy$. The multiplication is associative and commutative
if and only if $L_xL_y=L_yL_x$ for all $x,y\in M$.
\end{lemma}
\begin{proof}
Let $x,y\in M$. 
For an associative and commutative multiplication,
$L_xL_y=L_{xy}=L_{yx}=L_yL_x$. Conversely, applying $L_xL_y=L_yL_x$ to $1$
gives $xy=yx$. Hence, for all $z,t,u \in M$, we have 
\[
 (tu)z=z(tu)=L_zL_t(u)=L_tL_z(u)=t(zu)=t(uz)
\]
and the multiplication is associative as well.
\end{proof}

Recall that every fusion ring of rank four is commutative. 
Let $R$ be a free $\Z$-module of rank four with basis $\{b_0=1, b_1,b_2,b_3\}$ and multiplication $b_ib_j=\sum_k N_{ij}^k b_k$, where $N_{ij}^k \geq 0$ are integers, satisfying the commutativity, unit, duality, and Frobenius reciprocity axioms. 
Up to isomorphism, we can assume that either $b_i^*=b_i$ for all $i$, or $b_1^*=b_1$ and $b_2^*=b_3$. 

For every $u\in\{0,1,2,3\}$, let $N_u$ be the multiplication matrix of $b_u$. 
The matrix of the operator $L_{b_u}$, in the usual column convention, is $N_u^T$. 
Consequently, thanks to Lemma \ref{lem:commuting-multiplication}, $R$ is a fusion ring if and only if the matrices $N_1, N_2$, and $N_3$ pairwise commute.  

In the self-dual case, the multiplication matrices are of the form
\[
N_1=\begin{pmatrix}0&1&0&0\\1&a&b&c\\0&b&d&e\\0&c&e&f\end{pmatrix},\quad
N_2=\begin{pmatrix}0&0&1&0\\0&b&d&e\\1&d&g&h\\0&e&h&i\end{pmatrix},\quad
N_3=\begin{pmatrix}0&0&0&1\\0&c&e&f\\0&e&h&i\\1&f&i&j\end{pmatrix} 
\]
and their pairwise commutativity is equivalent to
\begin{align}
 ad-b^2+bg+ch-d^2-e^2+1&=0,\notag\\
 ae-bc+bh+ci-de-ef&=0,\notag\\
 be-cd+dh-eg+ei-fh&=0,\label{eq:sd-assoc}\\
 af+bi-c^2+cj-e^2-f^2+1&=0,\notag\\
 bf-ce+di-eh+ej-fi&=0,\notag\\
 df-e^2+gi-h^2+hj-i^2+1&=0.\notag
\end{align}
In the non-self-dual case, they are of the form \[
 N_1=\begin{pmatrix}0&1&0&0\\1&a&b&b\\0&b&c&d\\0&b&d&c\end{pmatrix},\quad
 N_2=\begin{pmatrix}0&0&1&0\\0&b&c&d\\0&d&e&f\\1&c&e&e\end{pmatrix},\quad
 N_3=N_2^{\mathsf T}
\]
and their pairwise commutativity is equivalent to
\begin{align}
 ac-b^2+2be-c^2-d^2+1&=0,\notag\\
 ad-b^2+be+bf-2cd&=0,\label{eq:ns-assoc}\\
 b(c-d)+d(e-f)&=0,\notag\\
 -c^2+d^2-e^2+f^2-1&=0.\notag
\end{align}

The following commands compute all solutions of \eqref{eq:sd-assoc} or \eqref{eq:ns-assoc} with $\FPdim(R)<3600$ (and so $\mu(R)\leq 59$ by Lemma \ref{lem:commutative}); they are $8977$ up to fusion ring isomorphism:
\begin{verbatim}
python3 supplementary/census/generate.py   
python3 supplementary/census/verify_census.py 
\end{verbatim}

Run all commands from the package root, alongside
\texttt{manuscript.tex} and the \texttt{supplementary} directory.
They require Python~3.10 or later and SymPy; census generation also
requires a C++17 compiler. The supplementary guide gives installation
instructions. No external census is used.

Theorem \ref{prop:subring-Fibo} excludes $1676$ rings containing a rank-two fusion subring with $x^2=1+ax$ and $a\geq2$.
The command
\begin{verbatim}
python3 supplementary/stages.py rank-two
\end{verbatim}
checks these exclusions and leaves $7301$ rings.
Among these rings, the command
\begin{verbatim}
python3 supplementary/stages.py d-number
\end{verbatim}
rules out $6541$ fusion rings that have a formal codegree failing the
$d$-number condition of Theorem \ref{prop:d-number}. 
It remains $760$ fusion rings.

Let $R$ be one of these $760$ fusion rings and assume that it is the Grothendieck ring of some spherical fusion category $\mathcal{C}$. 
Each character of $R$ gives candidate values for $\dim_\mathcal{C}(V)$ for every simple object $V\in \mathcal{O}(\mathcal{C})$. 
These numbers must satisfy some conditions according to Subsections \ref{sub-sec2.2} and \ref{sub-sec2.4}: they are nonzero real numbers, $\dim_\mathcal{C}(V)=\dim_\mathcal{C}(V^*)$ for all simple objects $V$, Ostrik's inequality \eqref{eq:spherical-mass}, and, for every formal codegree $f$ of $R$, the field $\Q(\dim(\mathcal{C})/f)$ is an abelian extension of $\Q$ and $\dim(\mathcal{C})/f$ has integral norm. 
The command
\begin{verbatim}
python3 supplementary/stages.py spherical
\end{verbatim}
eliminates $595$ rings for which no character satisfies all of the above conditions, leaving $165$ rings.


Among these $165$ rings, the published family restrictions in
Subsection~\ref{sub-sec2.4} give $90$ further exclusions. The command
\begin{verbatim}
python3 supplementary/stages.py families
\end{verbatim}
checks the full fusion rules of those families and leaves $75$ rings.

Now, if $R$ is one of these $75$ fusion rings, we write $S_R$ for the set of characters of $R$ passing all the arithmetic conditions above. 
This is a non-empty set by construction. 
Among these $75$ fusion rings $R$, we consider those for which $\dim(\mathcal{C})$ and $\FPdim(\mathcal{C})$ are Galois conjugates for all characters in $S_R$. 
For each character in $S_R$, Lemma \ref{lem:galois-reduction} tells us that $\mathcal{C}$ can be assumed pseudo-unitary, possibly after replacing $\mathcal{C}$ by a Galois conjugate depending on the character.
The command 
\begin{verbatim}
python3 supplementary/stages.py central
\end{verbatim}
excludes $60$ rings in this subset: every possible low induction matrix violates at least one necessary condition in Section~\ref{sec:central}.
There remain $15$ fusion rings, identified with Table~\ref{tab:models} by the command
\begin{verbatim}
python3 supplementary/stages.py models
\end{verbatim}
This proves Theorem~\ref{thm:main}. \qed


\par\smallskip
The non-necessarily pivotal case is the subject of our work
in progress, \emph{Classifying Grothendieck rings of general fusion
categories of rank four}.

\subsection*{Acknowledgments}
S\'ebastien Palcoux is supported by the National Natural Science
Foundation of China (NSFC), Grant no.~12471031.
Arnaud Plessis is supported by the Beijing Municipal Natural Science
Foundation (No.~IS25018).
GPT-6 Astra was used as a research assistant. The authors take full
responsibility for the final manuscript.

\raggedbottom
\begin{thebibliography}{99}

\bibitem{IntegralRings}
M. A. Alekseyev, W. Bruns, J. Dong and S. Palcoux,
\emph{Classifying integral Grothendieck rings up to rank 5 and beyond},
Lett. Math. Phys. \textbf{116} (2026), article 141.
\href{https://doi.org/10.1007/s11005-026-02164-7}{doi:10.1007/s11005-026-02164-7}.


\bibitem{BarrettWestbury}
J. W. Barrett and B. W. Westbury,
\emph{Invariants of piecewise-linear $3$-manifolds},
Trans. Amer. Math. Soc. \textbf{348} (1996), 3997--4022.
\href{https://doi.org/10.1090/S0002-9947-96-01660-1}{doi:10.1090/S0002-9947-96-01660-1}.

\bibitem{Bruillard4}
P. Bruillard,
\emph{Rank 4 premodular categories}, with an appendix by C. Galindo,
S.-H. Ng, J. Plavnik, E. Rowell and Z. Wang,
New York J. Math. \textbf{22} (2016), 775--800.
\href{https://nyjm.albany.edu/j/2016/22-36.html}{Journal article}.

\bibitem{DongSelfDual4}
J. Dong,
\emph{On Grothedieck rings of rank $4$ self-dual fusion categories},
\href{https://arxiv.org/abs/2505.19153v2}{arXiv:2505.19153v2} (2025).

\bibitem{DongChenWang}
J. Dong, G. Chen and Z. Wang,
\emph{Fusion categories containing a fusion subcategory with maximal rank},
J. Algebra \textbf{604} (2022), 107--127.
\href{https://doi.org/10.1016/j.jalgebra.2022.02.020}{doi:10.1016/j.jalgebra.2022.02.020}.

\bibitem{DongZhangDai}
J. Dong, L. Zhang and L. Dai,
\emph{Non-trivially graded self-dual fusion categories of rank $4$},
Acta Math. Sin. (Engl. Ser.) \textbf{34} (2018), 275--287.
\href{https://doi.org/10.1007/s10114-017-6375-0}{doi:10.1007/s10114-017-6375-0}.

\bibitem{EIP}
C. Edie-Michell, M. Izumi and D. Penneys,
\emph{Classification of $\mathbb Z/2\mathbb Z$-quadratic unitary fusion categories},
with an appendix by R. Johnson, S.-H. Ng, D. Penneys, J. Roat, M. Titsworth and H. Tucker,
Quantum Topol. \textbf{16} (2025), 493--568.
\href{https://doi.org/10.4171/QT/201}{doi:10.4171/QT/201}.

\bibitem{EGNO}
P. Etingof, S. Gelaki, D. Nikshych and V. Ostrik,
\emph{Tensor categories},
Mathematical Surveys and Monographs \textbf{205},
Amer. Math. Soc., Providence, RI, 2015.
\href{https://doi.org/10.1090/surv/205}{doi:10.1090/surv/205}.

\bibitem{ENO}
P. Etingof, D. Nikshych, and V. Ostrik,
\emph{On fusion categories},
Ann. of Math. (2) \textbf{162} (2005), no.~2, 581--642.
\href{https://doi.org/10.4007/annals.2005.162.581}{doi:10.4007/annals.2005.162.581}.

\bibitem{EvansJones}
D. E. Evans and C. Jones,
\emph{Quantum symmetries of noncommutative tori},
Comm. Math. Phys. \textbf{406} (2025), article 167.
\href{https://doi.org/10.1007/s00220-025-05337-y}{doi:10.1007/s00220-025-05337-y}.

\bibitem{GrossmanSnyder}
P. Grossman and N. Snyder,
\emph{Quantum subgroups of the Haagerup fusion categories},
Comm. Math. Phys. \textbf{311} (2012), 617--643.
\href{https://doi.org/10.1007/s00220-012-1427-x}{doi:10.1007/s00220-012-1427-x}.

\bibitem{HaggeHong}
T. J. Hagge and S.-M. Hong,
\emph{Some non-braided fusion categories of rank three},
Commun. Contemp. Math. \textbf{11} (2009), no.~4, 615--637.
\href{https://doi.org/10.1142/S0219199709003521}{doi:10.1142/S0219199709003521}.

\bibitem{HongRowell}
S.-M. Hong and E. C. Rowell,
\emph{On the classification of the Grothendieck rings of non-self-dual modular categories},
J. Algebra \textbf{324} (2010), 1000--1015.
\href{https://doi.org/10.1016/j.jalgebra.2009.11.044}{doi:10.1016/j.jalgebra.2009.11.044}.

\bibitem{HornJohnson}
R. A. Horn and C. R. Johnson,
\emph{Matrix analysis}, second edition,
Cambridge University Press, Cambridge, 2013.
\href{https://doi.org/10.1017/CBO9781139020411}{doi:10.1017/CBO9781139020411}.

\bibitem{IzumiE8}
M. Izumi,
\emph{On flatness of the Coxeter graph $E_8$},
Pacific J. Math. \textbf{166} (1994), 305--327.
\href{https://doi.org/10.2140/pjm.1994.166.305}{doi:10.2140/pjm.1994.166.305}.

\bibitem{IzumiNear}
M. Izumi,
\emph{A Cuntz algebra approach to the classification of near-group categories},
in \emph{Proceedings of the 2014 Maui and 2015 Qinhuangdao Conferences in Honour of
Vaughan F. R. Jones' 60th Birthday},
Proc. Centre Math. Appl. Austral. Nat. Univ. \textbf{46} (2017), 222--343.
\href{https://maths.anu.edu.au/files/8\%2520Izumi_2.pdf}{Published chapter}.

\bibitem{Larson}
H. K. Larson,
\emph{Pseudo-unitary non-self-dual fusion categories of rank 4},
J. Algebra \textbf{415} (2014), 184--213.
\href{https://doi.org/10.1016/j.jalgebra.2014.05.032}{doi:10.1016/j.jalgebra.2014.05.032}.

\bibitem{LevinWen}
M. A. Levin and X.-G. Wen,
\emph{String-net condensation: A physical mechanism for topological phases},
Phys. Rev. B \textbf{71} (2005), article 045110.
\href{https://doi.org/10.1103/PhysRevB.71.045110}{doi:10.1103/PhysRevB.71.045110}.

\bibitem{LMP}
Z. Liu, S. Morrison and D. Penneys,
\emph{1-supertransitive subfactors with index at most $6\frac15$},
Comm. Math. Phys. \textbf{334} (2015), 889--922.
\href{https://doi.org/10.1007/s00220-014-2160-4}{doi:10.1007/s00220-014-2160-4}.

\bibitem{LPR}
Z. Liu, S. Palcoux and Y. Ren,
\emph{Classification of Grothendieck rings of complex fusion categories of multiplicity one up to rank six},
Lett. Math. Phys. \textbf{112} (2022), article 54.
\href{https://doi.org/10.1007/s11005-022-01542-1}{doi:10.1007/s11005-022-01542-1}.

\bibitem{LPR2}
Z. Liu, S. Palcoux and Y. Ren,
\emph{Interpolated family of non-group-like simple integral fusion rings of {L}ie type},
Int. J. Math. \textbf{34} (2023), no.6, Paper No. 2350030, 51 pp.
\href{https://doi.org/10.1142/S0129167X23500301}{doi:10.1142/S0129167X23500301}.

\bibitem{NgSchauenburgHigher}
S.-H. Ng and P. Schauenburg,
\emph{Higher Frobenius--Schur indicators for pivotal categories},
in \emph{Hopf algebras and generalizations},
Contemp. Math. \textbf{441}, Amer. Math. Soc., Providence, RI, 2007, 63--90.
\href{https://doi.org/10.1090/conm/441/08500}{doi:10.1090/conm/441/08500}.

\bibitem{NgSchauenburg}
S.-H. Ng and P. Schauenburg,
\emph{Frobenius--Schur indicators and exponents of spherical categories},
Adv. Math. \textbf{211} (2007), 34--71.
\href{https://doi.org/10.1016/j.aim.2006.07.017}{doi:10.1016/j.aim.2006.07.017}.

\bibitem{OstrikRank2}
V. Ostrik,
\emph{Fusion categories of rank 2},
Math. Res. Lett. \textbf{10} (2003), 177--183.
\href{https://doi.org/10.4310/MRL.2003.v10.n2.a5}{doi:10.4310/MRL.2003.v10.n2.a5}.

\bibitem{OstrikCodegrees}
V. Ostrik,
\emph{On formal codegrees of fusion categories},
Math. Res. Lett. \textbf{16} (2009), 895--901.
\href{https://doi.org/10.4310/MRL.2009.v16.n5.a11}{doi:10.4310/\allowbreak MRL.2009.v16.n5.a11}.

\bibitem{OstrikRank3}
V. Ostrik,
\emph{Pivotal fusion categories of rank 3},
Mosc. Math. J. \textbf{15} (2015), 373--396.
\href{https://doi.org/10.17323/1609-4514-2015-15-2-373-396}{doi:10.17323/1609-4514-2015-15-2-373-396}.

\bibitem{OstrikGlobal}
V. Ostrik,
\emph{Remarks on global dimensions of fusion categories},
with an appendix joint with P. Etingof,
in \emph{Tensor Categories and Hopf Algebras},
Contemp. Math. \textbf{728}, Amer. Math. Soc., Providence, RI, 2019,
169--180.
\href{https://doi.org/10.1090/conm/728/14661}{doi:10.1090/conm/728/14661}.

\bibitem{RSW}
E. Rowell, R. Stong and Z. Wang,
\emph{On classification of modular tensor categories},
Comm. Math. Phys. \textbf{292} (2009), 343--389.
\href{https://doi.org/10.1007/s00220-009-0908-z}{doi:10.1007/s00220-009-0908-z}.

\bibitem{Schopieray}
A. Schopieray,
\emph{Categorification of integral group rings extended by one dimension},
J. Lond. Math. Soc. (2) \textbf{108} (2023), 1617--1641.
\href{https://doi.org/10.1112/jlms.12789}{doi:10.1112/jlms.12789}.

\bibitem{TambaraYamagami}
D. Tambara and S. Yamagami,
\emph{Tensor categories with fusion rules of self-duality for finite abelian groups},
J. Algebra \textbf{209} (1998), 692--707.
\href{https://doi.org/10.1006/jabr.1998.7558}{doi:10.1006/jabr.1998.7558}.

\end{thebibliography}

\end{document}
